% GENERATED FILE. Edit canonical lessons, then run npm run generate:books.
% canonical-source-hash: fnv1a64:fd488cc5c2a17bf9
% canonical-lessons: ZH-C433-jiedai, ZH-C433-haiwai, ZH-W433-hai-sea, ZH-C433-renshanrenhai, ZH-C433-huaqiao, ZH-C433-chuhai

\chapter{From Overseas}
\label{ch:zh-from-overseas}
\hlchaptermodality{433}

\begin{chapteropening}
\textbf{By the end of this chapter:} \emph{I can talk about receiving visitors from overseas, huge crowds, Chinese people abroad and going out to sea.}

Complete the last lesson of chapter 433: I can talk about receiving visitors from overseas, huge crowds, Chinese people abroad and going out to sea.
\end{chapteropening}

\section[jiēdài]{jiēdài — to receive visitors; to take care of guests}

\label{lesson:ZH-C433-jiedai}



\begin{quote}
Before the new one: say the Mandarin for a full-price ticket, then the Mandarin for to be shut to the public.
\end{quote}

\subsection*{You'll want to know: jiēdài}
\emph{jiēdài} — \textquotedblleft{}to receive visitors; to take care of guests\textquotedblright{}.

\begin{quote}
\emph{zhè jiā jiǔdiàn jiēdài le hěn duō wàiguó yóukè} — This hotel has had lots of foreign visitors.
\end{quote}

Its characters come later; for now it is a word for the ear.

\begin{cousinweb}
\emph{jiēdài} is built from \emph{jiē} 'to meet' and \emph{dài} 'to wait on', both inherited from Middle Chinese, \emph{tsjep} and \emph{dojX} (Guangyun, in Baxter's transcription), which glosses the first 'to join' and the second 'to wait for'.
\end{cousinweb}

\begin{grammarlens}[title={What you've built}]
One more word for getting around.
\end{grammarlens}

\subsection*{Your turn}
\begin{itemize}
\raggedright
  \item \emph{Say it:} \emph{jiēdài}
  \item \emph{Say it:} \emph{jiēdài}, once more
  \item \emph{Say it:} \emph{zhè jiā jiǔdiàn jiēdài le hěn duō wàiguó yóukè} — This hotel has had lots of foreign visitors.
  \item \emph{Recall:} say the Mandarin for a full-price ticket, then the Mandarin for to be shut to the public, then say \emph{jiēdài} again
\end{itemize}

\begin{tcolorbox}[breakable,colback=teal!4,colframe=teal!35!black,title={Before you move on}]
What does jiēdài mean? (\textquotedblleft{}To receive visitors; to take care of guests\textquotedblright{}.) Say it once more.
\end{tcolorbox}
\section[\texorpdfstring{\zh{海外} (hǎiwài)}{海外 (hǎiwài)}]{\zh{海外} (hǎiwài) — overseas; beyond the sea}

\label{lesson:ZH-C433-haiwai}



\begin{quote}
Before the new one: say the Mandarin for to be shut to the public, then the Mandarin for to receive visitors.
\end{quote}

\subsection*{You'll want to know: \zh{海外}}
\textbf{\zh{海外}} — \emph{hǎiwài} — \textquotedblleft{}overseas; beyond the sea\textquotedblright{}.

\begin{quote}
\emph{hěn duō hǎiwài yóukè zhù zài zhè jiā jiǔdiàn} — A lot of visitors from overseas stay at this hotel.
\end{quote}

The new character, in pieces: \zh{海} — ten strokes in two pieces: on the left the water piece of \zh{洗}, three strokes; on the right \zh{每} měi, seven strokes. The next lesson puts it in your hand.

\begin{cousinweb}
\textbf{\zh{海外}} (\emph{hǎiwài}) is built from \zh{海} 'the sea' and \zh{外} 'outside', both inherited from Middle Chinese, \emph{xojX} and \emph{ngwajH} (Guangyun, in Baxter's transcription), which glosses the first, after an older dictionary, as the great pool that takes in a hundred rivers, and the second 'outside'.
\end{cousinweb}

\begin{grammarlens}[title={What you've built}]
One more word for getting around.
\end{grammarlens}

\subsection*{Your turn}
\begin{itemize}
\raggedright
  \item \emph{Say it:} \emph{hǎiwài}
  \item \emph{Say it:} \emph{hǎiwài}, once more
  \item \emph{Say it:} \emph{hěn duō hǎiwài yóukè zhù zài zhè jiā jiǔdiàn} — A lot of visitors from overseas stay at this hotel.
  \item \emph{Recall:} say the Mandarin for to be shut to the public, then the Mandarin for to receive visitors, then say \emph{hǎiwài} again
\end{itemize}

\begin{tcolorbox}[breakable,colback=teal!4,colframe=teal!35!black,title={Before you move on}]
What does \zh{海外} mean? (\textquotedblleft{}Overseas; beyond the sea\textquotedblright{}.) Say it once more.
\end{tcolorbox}
\section[\texorpdfstring{\zh{海} (hǎi)}{海 (hǎi)}]{\zh{海} — the water piece, then \zh{每}}

\label{lesson:ZH-W433-hai-sea}



\begin{quote}
The word first, then the character inside it.

\begin{itemize}
\raggedright
  \item \emph{Recall:} say \emph{hǎiwài} — \textbf{R1}, one lesson back, before it decays
\end{itemize}

\begin{quote}
\zh{洗} \zh{每}
\end{quote}

\emph{Read it:} each one aloud, with nothing said first — \emph{xǐ}, \emph{měi}
\end{quote}

\subsection*{Script — the water piece, then \zh{每}}
\hlblockfigure{\detokenize{figures/ZH-W433-hai-sea-filmstrip.pdf}}{How \zh{海} is written, stroke by stroke}

\begin{quote}
\zh{海}
\end{quote}

\textbf{Two pieces, left before right.} Count the pieces, not the strokes.

\textbf{The left piece, the water piece, three strokes}, as in \zh{洗}: two dots, then a stroke that rises to the right.

\textbf{The right piece, \zh{每} měi, seven strokes}, as you wrote it: a short stroke falling to the left and a bar; a stroke that runs down and turns right along the foot; a bar that turns down and hooks; a dot; a long bar across; and a last dot.

3 + 7 = ten.

\subsection*{Your turn}
\begin{itemize}
\raggedright
  \item \emph{Write it:} \zh{海}, all ten, with the model in front of you
  \item \emph{Cover:} write it again from memory, and count the pieces, not the strokes
\end{itemize}

\begin{tcolorbox}[breakable,colback=teal!4,colframe=teal!35!black,title={Before you move on}]
How many strokes in \zh{海}? (\textbf{Ten}.) What is the right piece? (\textbf{\zh{每} měi}.) How many dots does \zh{每} have? (\textbf{Two}.)

Source: \href{https://raw.githubusercontent.com/chanind/hanzi-writer-data/68d10a4b21150cae5e1ebbd223eed289cf32d90c/data/\%E6\%B5\%B7.json}{Hanzi Writer Data \zh{海}, pinned stroke paths and medians}.
\end{tcolorbox}
\section[\texorpdfstring{\zh{人山人海} (rénshānrénhǎi)}{人山人海 (rénshānrénhǎi)}]{\zh{人山人海} (rénshānrénhǎi) — huge crowds; a sea of people}

\label{lesson:ZH-C433-renshanrenhai}



\begin{quote}
Before the new one: say the Mandarin for to receive visitors, then the Mandarin for beyond the sea.
\end{quote}

\subsection*{You'll want to know: \zh{人山人海}}
\textbf{\zh{人山人海}} — \emph{rénshānrénhǎi} — \textquotedblleft{}huge crowds; a sea of people\textquotedblright{}.

\begin{quote}
\emph{zhōumò gōngyuán lǐ rénshānrénhǎi} — At the weekend the park is packed with people.
\end{quote}

\begin{cousinweb}
\textbf{\zh{人山人海}} (\emph{rénshānrénhǎi}) is built from \zh{人} 'people', \zh{山} 'a mountain' and \zh{海} 'the sea', said twice with \zh{人}: people like mountains, people like seas. The parts are inherited from Middle Chinese, \emph{nyin}, \emph{srean} and \emph{xojX} (Guangyun, in Baxter's transcription).
\end{cousinweb}

\begin{grammarlens}[title={What you've built}]
One more word for getting around.
\end{grammarlens}

\subsection*{Your turn}
\begin{itemize}
\raggedright
  \item \emph{Say it:} \emph{rénshānrénhǎi}
  \item \emph{Say it:} \emph{rénshānrénhǎi}, once more
  \item \emph{Say it:} \emph{zhōumò gōngyuán lǐ rénshānrénhǎi} — At the weekend the park is packed with people.
  \item \emph{Recall:} say the Mandarin for to receive visitors, then the Mandarin for beyond the sea, then say \emph{rénshānrénhǎi} again
\end{itemize}

\begin{tcolorbox}[breakable,colback=teal!4,colframe=teal!35!black,title={Before you move on}]
What does \zh{人山人海} mean? (\textquotedblleft{}Huge crowds; a sea of people\textquotedblright{}.) Say it once more.
\end{tcolorbox}
\section[huáqiáo]{huáqiáo — overseas Chinese; a Chinese person settled abroad}

\label{lesson:ZH-C433-huaqiao}



\begin{quote}
Before the new one: say the Mandarin for beyond the sea, then the Mandarin for a sea of people.
\end{quote}

\subsection*{You'll want to know: huáqiáo}
\emph{huáqiáo} — \textquotedblleft{}overseas Chinese; a Chinese person settled abroad\textquotedblright{}.

\begin{quote}
\emph{tā shì zhù zài hǎiwài de huáqiáo} — He's one of the Chinese who live overseas.
\end{quote}

Its characters come later; for now it is a word for the ear.

\begin{cousinweb}
\emph{huáqiáo} is built from \emph{huá} 'Chinese' and \emph{qiáo} 'to live away from home', both inherited from Middle Chinese, \emph{hwae} and \emph{gjew} (Guangyun, in Baxter's transcription), which glosses the first 'flourishing' and the second 'to lodge' and 'a guest'; the sense 'Chinese' for the first is not given there.
\end{cousinweb}

\begin{grammarlens}[title={What you've built}]
One more word for getting around.
\end{grammarlens}

\subsection*{Your turn}
\begin{itemize}
\raggedright
  \item \emph{Say it:} \emph{huáqiáo}
  \item \emph{Say it:} \emph{huáqiáo}, once more
  \item \emph{Say it:} \emph{tā shì zhù zài hǎiwài de huáqiáo} — He's one of the Chinese who live overseas.
  \item \emph{Recall:} say the Mandarin for beyond the sea, then the Mandarin for a sea of people, then say \emph{huáqiáo} again
\end{itemize}

\begin{tcolorbox}[breakable,colback=teal!4,colframe=teal!35!black,title={Before you move on}]
What does huáqiáo mean? (\textquotedblleft{}Overseas Chinese; a Chinese person settled abroad\textquotedblright{}.) Say it once more.
\end{tcolorbox}
\section[\texorpdfstring{\zh{出海} (chūhǎi)}{出海 (chūhǎi)}]{\zh{出海} (chūhǎi) — to go out to sea}

\label{lesson:ZH-C433-chuhai}



\begin{quote}
Before the new one: say the Mandarin for a sea of people, then the Mandarin for a Chinese person settled abroad.
\end{quote}

\subsection*{You'll want to know: \zh{出海}}
\textbf{\zh{出海}} — \emph{chūhǎi} — \textquotedblleft{}to go out to sea\textquotedblright{}.

\begin{quote}
\textbf{\zh{今天风很大我们不能出海}} \emph{jīntiān fēng hěn dà wǒmen bù néng chūhǎi} — It's very windy today, so we can't go out to sea.
\end{quote}

\begin{cousinweb}
\textbf{\zh{出海}} (\emph{chūhǎi}) is built from \zh{出} 'to go out' and \zh{海} 'the sea', both inherited from Middle Chinese, \emph{tsyhwit} and \emph{xojX} (Guangyun, in Baxter's transcription).
\end{cousinweb}

\begin{grammarlens}[title={What you've built}]
One more word for getting around.
\end{grammarlens}

\subsection*{Your turn}
\begin{itemize}
\raggedright
  \item \emph{Say it:} \emph{chūhǎi}
  \item \emph{Say it:} \emph{chūhǎi}, once more
  \item \emph{Say it:} \emph{jīntiān fēng hěn dà wǒmen bù néng chūhǎi} — It's very windy today, so we can't go out to sea.
  \item \emph{Recall:} say the Mandarin for a sea of people, then the Mandarin for a Chinese person settled abroad, then say \emph{chūhǎi} again
\end{itemize}

\begin{tcolorbox}[breakable,colback=teal!4,colframe=teal!35!black,title={Before you move on}]
What does \zh{出海} mean? (\textquotedblleft{}To go out to sea\textquotedblright{}.) Say it once more.
\end{tcolorbox}
